\section{Shift-And Algorithm}
\label{sec:shift-and}

We turn to a different perspective on exact pattern matching---one that works
\emph{from the bottom up}, at the level of individual bits inside machine words.

\aside{The same matrix will reappear when we tackle
\emph{approximate} string matching (matching with errors).}

The Shift-And algorithm combines two ideas:
\begin{enumerate}[nosep]
  \item \textbf{Bit vectors.}  Represent states as sequences of bits (words of
        memory) so that many comparisons happen in a single CPU instruction.
  \item \textbf{Dynamic programming.}  Define a matrix that records which
        prefixes of~$P$ are ``alive'' at every position of~$T$, then compute it
        column by column.
\end{enumerate}

Before diving into the algorithm itself, let us briefly recall these two
ingredients.

% ── Intermezzo: bit vectors ──────────────────────────────
\begin{intermezzo}{Bit vectors and bit-level parallelism}
A \emph{bit vector} is simply a sequence of bits---zeros and ones---packed into
a machine word.  A modern CPU has registers of width $w$~bits (typically
$w = 64$); a single register can therefore store 64 independent boolean values.

The key advantage is \emph{parallelism for free}: a bitwise AND (\texttt{\&}),
OR (\texttt{\textbar}), or shift (\texttt{<{}<}, \texttt{>{}>}) on a $w$-bit register
operates on all~$w$ bits in one clock cycle.  If we encode a problem so that
each ``slot'' of the answer is one bit, we effectively perform $w$~operations at
the cost of one.

\smallskip\noindent
\textbf{Example.}  Let $A = 10110$ and $B = 11010$.  Then:
\begin{center}
\begin{tabular}{rl}
  $A \;\texttt{\&}\; B$ & $= 10010$ \quad (AND: both bits must be 1) \\
  $A \;\texttt{|}\; B$  & $= 11110$ \quad (OR: at least one bit is 1) \\
  $A \;\texttt{<{}<}\; 1$ & $= 01100$ \quad (left-shift by one position)
\end{tabular}
\end{center}
Each of these is a \emph{single CPU instruction}, regardless of how many bits
are in the word.
\end{intermezzo}

% ── Intermezzo: dynamic programming ────────────────────────
\begin{intermezzo}{Dynamic programming}
\emph{Dynamic programming} (often abbreviated DP) is a general technique for
solving problems that can be decomposed into overlapping subproblems.  The idea
is straightforward:

\begin{enumerate}[nosep]
  \item Define a table (matrix, array, \ldots) whose entries represent solutions
        to ``smaller'' instances of the problem.
  \item Write a \emph{recurrence}: a formula that computes each entry from
        entries that have already been filled in.
  \item Fill the table in an order that respects the dependencies, so that every
        entry we need is ready when we need it.
\end{enumerate}

The payoff is that we solve each subproblem exactly once and store its answer,
rather than recomputing it every time it is needed (which is what a naive
recursive approach would do).

In the context of the Shift-And algorithm, the ``table'' is a bit matrix~$M$
where $M[i,j]$ tells us whether a certain prefix of~$P$ matches a certain
suffix of~$T$.  We will fill this table column by column, left to right,
using the recurrence derived below.

\smallskip\noindent
\textbf{Important.}  The full $m \times n$ matrix is \emph{never materialised
in memory}.  Since each column depends only on the immediately preceding
column, we only need to keep \emph{one column at a time}---a single machine
word when $m \le w$.  This is a common space optimisation in dynamic programming.
\end{intermezzo}

% ── Section ──────────────────────────────────────────────
\subsection{Working assumption}

\begin{remark}
  We assume $|P| \le w$, where $w$ is the machine word size (typically~64).
  Under this assumption every column of the matrix fits in a single register,
  and shift/AND operations cost $\Oh{1}$.
\end{remark}

Working with a very large text~$T$ is not a problem: the algorithm scans~$T$
once, left to right, spending constant time per character.
When $|P| > w$ the bit vector is split into $\lceil m/w \rceil$ machine words
and the cost per text character rises to $\Oh{\lceil m/w \rceil}$.

% ── Section ──────────────────────────────────────────────
\subsection{The matrix}

Define a bit matrix $M$ of size $m \times n$ (rows indexed by the pattern,
columns by the text):

\begin{definition}[Shift-And matrix]
\label{def:sa-matrix}
  \[
    M[i,j] \;=\;
    \begin{cases}
      1 & \text{if } P[1 \mathinner{.\,.} i]
          \text{ is a suffix of } T[1 \mathinner{.\,.} j],\\[2pt]
      0 & \text{otherwise.}
    \end{cases}
  \]
\end{definition}

In words: $M[i,j] = 1$ exactly when the first $i$ characters of the pattern
match the last $i$ characters of $T[1 \mathinner{.\,.} j]$.

\begin{observation}
  An occurrence of~$P$ ending at position~$j$ of~$T$ corresponds to
  $M[m,j] = 1$.  Therefore, the set of all 1s in the \emph{last row} gives us
  every occurrence.
\end{observation}

\aside{The 1s in the other rows carry information too:
they mark partial matches that may extend as we read more of~$T$.}

As noted above, this matrix is never stored in full: the algorithm computes it
one column at a time, overwriting the previous column at each step
(Figure~\ref{fig:sa-matrix}).

% ── Figure: DP matrix diagram ──────────────────────────────
\begin{figure}[H]
\centering
\begin{tikzpicture}[
  cell/.style={
    minimum width=7mm, minimum height=7mm,
    font=\ttfamily\small, inner sep=0pt,
    text centered, rounded corners=2pt,
  },
  Tcell/.style={cell, draw=blue!50!black, fill=blue!8},
  Pcell/.style={cell, draw=accent!50!black, fill=accent!8},
  hitcell/.style={cell, draw=green!55!black, fill=green!18, font=\small\bfseries},
  blankcell/.style={cell, draw=gray!25, fill=gray!4, text=gray!50},
]
\def\cw{0.82}

% ── column index j ───────────────────────────────────────
\node[font=\small\sffamily, text=gray] at ({2.5*\cw}, 1.9) {$j$};
\foreach \j [count=\k from 0] in {1,...,6}
  \node[font=\tiny\sffamily, text=gray!70] at ({\k*\cw}, 1.5) {\j};

% ── T row (just above the matrix, with moderate gap) ──────
\node[font=\small\sffamily, anchor=east] at (-0.55, 0.9) {$T$};
\foreach \j/\c in {1/a, 2/b, 3/c, 4/a, 5/b, 6/c}
  \node[Tcell] at ({(\j-1)*\cw}, 0.9) {\c};

% ── row index i (far left, clear of P boxes) ─────────────
\node[font=\small\sffamily, text=gray] at (-2.1, {-1*\cw}) {$i$};
\foreach \i [count=\k from 0] in {1,2,3}
  \node[font=\tiny\sffamily, text=gray!70] at (-1.7, {-\k*\cw}) {\i};

% ── P column (left of matrix, shifted right of i indices) ─
\node[font=\small\sffamily, anchor=east] at (-2.45, {-1*\cw}) {$P$};
\foreach \i/\c in {1/a, 2/b, 3/c}
  \node[Pcell] at (-1.2, {-(\i-1)*\cw}) {\c};

% ── Matrix cells ──────────────────────────────────────────
% Row i=1
\foreach \j/\v in {1/1, 2/0, 3/0, 4/1, 5/0, 6/0} {
  \ifnum\v=1
    \node[hitcell] at ({(\j-1)*\cw}, 0) {\v};
  \else
    \node[blankcell] at ({(\j-1)*\cw}, 0) {0};
  \fi
}
% Row i=2
\foreach \j/\v in {1/0, 2/1, 3/0, 4/0, 5/1, 6/0} {
  \ifnum\v=1
    \node[hitcell] at ({(\j-1)*\cw}, {-1*\cw}) {\v};
  \else
    \node[blankcell] at ({(\j-1)*\cw}, {-1*\cw}) {0};
  \fi
}
% Row i=3 (last row)
\foreach \j/\v in {1/0, 2/0, 3/1, 4/0, 5/0, 6/1} {
  \ifnum\v=1
    \node[hitcell] at ({(\j-1)*\cw}, {-2*\cw}) {\v};
  \else
    \node[blankcell] at ({(\j-1)*\cw}, {-2*\cw}) {0};
  \fi
}

% ── Highlight last row ────────────────────────────────────
\draw[green!55!black, thick, rounded corners=3pt]
  ([xshift=-5pt, yshift=-5pt] {0*\cw-0.35},{-2*\cw - 0.35})
  rectangle
  ([xshift=5pt, yshift=5pt] {5*\cw+0.35},{-2*\cw + 0.35});

\node[font=\footnotesize\sffamily, text=green!45!black, anchor=west]
  at ({5*\cw + 0.65}, {-2*\cw}) {occurrences};

\end{tikzpicture}
\caption{%
  Shift-And matrix for $P = \texttt{abc}$ and $T = \texttt{abcabc}$.
  The 1s in the last row ($i = m = 3$) mark positions $j = 3$ and $j = 6$
  where an occurrence of~$P$ ends.
  Note that the full matrix is never stored: the algorithm keeps only one
  column at a time.
}
\label{fig:sa-matrix}
\end{figure}

% ── Section ──────────────────────────────────────────────
\subsection{Recurrence}

How do we compute $M[i,j]$?  A prefix $P[1 \mathinner{.\,.} i]$ is a suffix of
$T[1 \mathinner{.\,.} j]$ if and only if two conditions hold simultaneously:

\begin{enumerate}[nosep]
  \item $P[1 \mathinner{.\,.} i{-}1]$ is already a suffix of
        $T[1 \mathinner{.\,.} j{-}1]$
        \aside{This is the \emph{shift}: the shorter prefix was alive at the
        previous position.}
        (i.e.\ $M[i{-}1,\, j{-}1] = 1$), and
  \item The last characters match: $P[i] = T[j]$.
\end{enumerate}

Formally:
\[
  M[i,j] = 1
  \quad\iff\quad
  M[i{-}1,\, j{-}1] = 1
  \;\;\text{AND}\;\;
  P[i] = T[j],
\]
with the convention that $M[0, j] = 1$ for all~$j$ (the empty prefix is always
a suffix).

This is where the name comes from: condition~(a) is a \textbf{shift} (look one
row up and one column left), and condition~(b) is an \textbf{AND} with the
alphabet mask.

% ── Section ──────────────────────────────────────────────
\subsection{Bit-vector reformulation}

The key insight is that we can compute an entire column of~$M$ in one shot by
encoding it as a bit vector.

\subsubsection*{Alphabet masks}

For each character $x \in \Sigma$, pre-compute a bit vector $U_x$ of length~$m$:
\[
  U_x[i] =
  \begin{cases}
    1 & \text{if } P[i] = x,\\
    0 & \text{otherwise.}
  \end{cases}
\]
This records every position in the pattern where character~$x$ appears.

\begin{example}
  For $P = \texttt{abab}$ over $\Sigma = \{a,b\}$:
  \[
    U_a = 1010, \qquad U_b = 0101.
  \]
\end{example}

\subsubsection*{Column update}

Let $M_j$ denote the $j$-th column of~$M$, viewed as a single bit vector of
length~$m$ (bit~1 is the top, bit~$m$ is the bottom).  The recurrence becomes:

\[
  M_j \;=\; \bigl(\operatorname{shift}(M_{j-1}) \;\mathbin{\texttt{\&}}\; U_{T[j]}\bigr),
\]
where $\operatorname{shift}$ means:
\begin{itemize}[nosep]
  \item Shift the vector \emph{down} by one position (toward higher indices),
  \item Insert a~$1$ at the top (position~1), because the empty prefix is
        always alive,
  \item The bit that falls off the bottom ($M[m]$) is discarded.
\end{itemize}

\aside{In C-like code the ``shift down'' on a bit vector stored
LSB\,=\,position~1 is simply a left shift: \texttt{(M\,<{}<\,1)\,|\,1}.}

On a machine register, the shift-and-mask operation
(illustrated step by step in Figure~\ref{fig:sa-step}) becomes:
\[
  M_j \;=\; \bigl((M_{j-1} \ll 1) \mathbin{|} 1\bigr) \;\mathbin{\texttt{\&}}\; U_{T[j]}.
\]

% ── Figure: shift-and step ────────────────────────────────
\begin{figure}[H]
\centering
\begin{tikzpicture}[
  box/.style={
    minimum width=8mm, minimum height=7mm,
    draw=gray!50, fill=accent!6,
    font=\ttfamily\small, inner sep=0pt,
    text centered, rounded corners=2pt,
  },
  label/.style={font=\small\sffamily, text=gray!70},
  arrow/.style={-{Stealth[length=5pt]}, thick, accent!60},
]
\def\bw{0.9}

% Column M[j-1]
\node[label] at (0, 1.4) {$M_{j-1}$};
\foreach \i/\v in {1/b_1, 2/b_2, 3/b_3, 4/b_4}
  \node[box] (A\i) at (0, {-(\i-1)*\bw}) {$\v$};

% Arrow
\draw[arrow] (1.3, {-1.5*\bw}) -- node[above, font=\footnotesize\sffamily] {shift} (2.5, {-1.5*\bw});

% Shifted + OR 1
\node[label] at (3.8, 1.4) {$\operatorname{shift}$};
\node[box, fill=green!12] (B1) at (3.8, 0) {$1$};
\foreach \i/\v in {2/b_1, 3/b_2, 4/b_3}
  \node[box] (B\i) at (3.8, {-(\i-1)*\bw}) {$\v$};
\node[font=\tiny\sffamily, text=gray!50, anchor=west] at (4.4, {-4*\bw+0.3})
  {$b_4$ discarded};

% Arrow
\draw[arrow] (5.1, {-1.5*\bw}) -- node[above, font=\footnotesize\sffamily] {AND} (6.3, {-1.5*\bw});

% Mask U_{T[j]}
\node[label] at (7.6, 1.4) {$U_{T[j]}$};
\foreach \i/\v in {1/u_1, 2/u_2, 3/u_3, 4/u_4}
  \node[box, fill=accentgold!15] (C\i) at (7.6, {-(\i-1)*\bw}) {$\v$};

% Arrow
\draw[arrow] (8.9, {-1.5*\bw}) -- node[above, font=\footnotesize\sffamily] {$=$} (9.8, {-1.5*\bw});

% Result M[j]
\node[label] at (10.8, 1.4) {$M_j$};
\foreach \i in {1,...,4}
  \node[box, fill=accent!12] (D\i) at (10.8, {-(\i-1)*\bw}) {$\cdot$};

\end{tikzpicture}
\caption{%
  One step of the Shift-And algorithm.  The previous column is shifted down
  (with a~$1$ inserted at the top), then AND-ed with the alphabet mask for the
  current text character.
}
\label{fig:sa-step}
\end{figure}

% ── Section ──────────────────────────────────────────────
\subsection{Algorithm}

\begin{algorithm}[H]
\caption{Shift-And exact matching}
\label{alg:shift-and}
\KwIn{Text $T[1 \mathinner{.\,.} n]$, pattern $P[1 \mathinner{.\,.} m]$
      with $m \le w$}
\KwOut{All positions $j$ where $P$ occurs in $T$}
\BlankLine
\tcp{Pre-compute alphabet masks}
\ForEach{$x \in \Sigma$}{
  $U_x \leftarrow \mathbf{0}$\;
  \For{$i \leftarrow 1$ \KwTo $m$}{
    \If{$P[i] = x$}{
      $U_x \leftarrow U_x \mathbin{|} (1 \ll (i{-}1))$\;
    }
  }
}
\BlankLine
\tcp{Scan text}
$M \leftarrow 0$\;
\For{$j \leftarrow 1$ \KwTo $n$}{
  $M \leftarrow \bigl((M \ll 1) \mathbin{|} 1\bigr) \;\mathbin{\texttt{\&}}\; U_{T[j]}$\;
  \If{$M \;\mathbin{\texttt{\&}}\; (1 \ll (m{-}1)) \neq 0$}{
    output ``occurrence ending at position $j$''\;
  }
}
\end{algorithm}

% ── Section ──────────────────────────────────────────────
\subsection{Complexity}

\begin{itemize}[nosep]
  \item \textbf{Time.}
    Pre-processing the alphabet masks takes $\Oh{|\Sigma| \cdot m}$.
    When $m \le w$, each iteration of the main loop costs $\Oh{1}$
    (a shift, an OR, and an AND---three CPU instructions).
    The loop runs $n$ times, so the total search time is $\Oh{n}$.

    If $m > w$, each column requires $\lceil m/w \rceil$ words, and the
    cost per text character rises to $\Oh{\lceil m/w \rceil}$, giving
    $\Oh{n \cdot \lceil m/w \rceil}$ overall.

  \item \textbf{Space.}
    $\Oh{|\Sigma| \cdot \lceil m/w \rceil}$ for the masks, plus a single
    register for the current column~$M$.
    The full $m \times n$ matrix is \emph{never materialised}: we only
    keep one column at a time.
\end{itemize}

\begin{remark}
  The Shift-And algorithm is particularly attractive in bioinformatics
  (DNA sequencing), where the alphabet is small ($|\Sigma| = 4$), patterns
  are short, and the text can be very long.
  The same dynamic programming framework extends naturally to
  \emph{approximate matching} (allowing a bounded number of mismatches,
  insertions, or deletions)---a topic we will revisit later.
\end{remark}
